% Zone „Das kennst du schon“ – aus bank/ableitung-und-aenderungsrate/zone.jsonl (zusammenbau v0.1)
\einheitenkopf[zone][Kennst du schon]{Das kennst du schon}
\setcounter{aufgabe}{0}

%% TODO Titel: Ich-kann-Satz fehlt in der Bank (Platzhalter: Kettenname) – Nr. 1
%% TODO Anweisung: ein Satz über der ersten Teilaufgabe fehlt in der Bank – Nr. 1
\begin{aufgabe}{Steigung einer Geraden als Quotient aus Höhenunterschied und waagerechtem Unterschied, Steigungsdreieck lesen und zeichnen, Geradengleichung aus Steigung und Punkt (Punkt-Steigungs-Form)}
\begin{teile}
\teil Eine Gerade geht durch $A(1 | 2)$ und $B(3 | 8)$. Berechne ihre Steigung. m = \leerfeld
\teil Lies die Steigung der Geraden $g$ am Steigungsdreieck ab und gib die Gleichung von $g$ an.

\begin{ksys}[xmin=-2,xmax=4,ymin=-3,ymax=7,ablesen]\gerade{2}{-1}{g}\steigungsdreieck{1}{1}{1}\end{ksys}
y = \leerfeld
\end{teile}
\end{aufgabe}

%% TODO Titel: Ich-kann-Satz fehlt in der Bank (Platzhalter: Kettenname) – Nr. 2
%% TODO Anweisung: ein Satz über der ersten Teilaufgabe fehlt in der Bank – Nr. 2
\begin{aufgabe}{Funktionswerte berechnen, auch mit e-Term und Bruch, und Werte an Graph und Tabelle ablesen}
\begin{teile}
\teil $f(x) = x^2 - 3x$. Berechne $f(4)$. f(4) = \leerfeld
\teil $h(t) = 5t \cdot \mathrm{e}^{-0{,}2t}$. Berechne $h(3)$ auf zwei Nachkommastellen. h(3) = \leerfeld
\end{teile}
\end{aufgabe}

%% TODO Titel: Ich-kann-Satz fehlt in der Bank (Platzhalter: Kettenname) – Nr. 3
%% TODO Anweisung: ein Satz über der ersten Teilaufgabe fehlt in der Bank – Nr. 3
\begin{aufgabe}{Einheiten und Zeiten: Uhrzeiten in Stunden seit einem Bezugspunkt und Dezimalstunden in Stunden und Minuten umrechnen, Einheiten von f und x zur Einheit „je“ zusammensetzen (Stück je Tag, Meter je Minute), prozentuale Abweichung als Quotient}
\begin{teile}
\teil Die Zählung beginnt um 8:00 Uhr. Wie viele Stunden sind um 11:30 Uhr vergangen? \leerfeld[h]
\teil Rechne $2{,}75$ Stunden in Stunden und Minuten um, und gib an, um wie viel Prozent $9{,}2$ von $8$ abweicht. \leerfeld[h] \leerfeld[min] ; \leerfeld[\%]
\end{teile}
\end{aufgabe}

%% TODO Titel: Ich-kann-Satz fehlt in der Bank (Platzhalter: Kettenname) – Nr. 4
%% TODO Anweisung: ein Satz über der ersten Teilaufgabe fehlt in der Bank – Nr. 4
\begin{aufgabe}{Ableitungsfunktion bilden: Potenz-, Faktor- und Summenregel für ganzrationale Funktionen; GK und LK zusätzlich Produkt- und Kettenregel für Produkte aus Polynom und e-Funktion}
\begin{teile}
\teil Bilde $f'(x)$ für $f(x) = x^3 - 4x$. f'(x) = \leerfeld
\teil Bilde $h'(t)$ für $h(t) = t \cdot \mathrm{e}^{-0{,}5t}$. h'(t) = \leerfeld
\end{teile}
\end{aufgabe}

%% TODO Titel: Ich-kann-Satz fehlt in der Bank (Platzhalter: Kettenname) – Nr. 5
%% TODO Anweisung: ein Satz über der ersten Teilaufgabe fehlt in der Bank – Nr. 5
\begin{aufgabe}{Gleichungen lösen: linear, quadratisch mit Lösungsformel, Ausklammern und Nullprodukt; GK zusätzlich e-Faktoren als nullstellenfrei erkennen und einfache Exponentialgleichungen logarithmieren; mit dem Rechner numerisch lösen, wenn Teil B es zulässt}
\begin{gleichungsraster}[2]
\gl{\text{Löse: $4t - 10 = 6$}} &
\gl{\text{Löse: $3t^2 - 18t = 0$}} \\
\end{gleichungsraster}
\end{aufgabe}

%% TODO Titel: Ich-kann-Satz fehlt in der Bank (Platzhalter: Kettenname) – Nr. 6
%% TODO Anweisung: ein Satz über der ersten Teilaufgabe fehlt in der Bank – Nr. 6
\begin{aufgabe}{Verlauf eines Graphen qualitativ lesen: steigt, fällt, wird steiler oder flacher, Wendestelle als steilste Stelle, Krümmung}
\begin{teile}
\teil Der Graph zeigt eine Funktion $f$. In welchem Bereich steigt der Graph? \\ \kreuz{links von $x = 1$} \\ \kreuz{rechts von $x = 1$} \\ \kreuz{überall}

\begin{ksys}[xmin=-2,xmax=4,ymin=-2,ymax=6,ablesen]\parabel{1}{1}{-1}{f}\end{ksys}
\teil Der Graph zeigt die Höhe eines Baums in Abhängigkeit von der Zeit. Beschreibe in einem Satz, wie sich das Wachstum im gezeigten Zeitraum verändert.

\begin{ksys}[xmin=0,xmax=8,ymin=0,ymax=8,xstep=1,ystep=1,xlabel=t in Jahren,ylabel=h in m,ablesen]\funktion{8/(1+exp(-1.2*(\x-4)))}{h}\end{ksys}
\end{teile}
\end{aufgabe}

%% TODO Titel: Ich-kann-Satz fehlt in der Bank (Platzhalter: Kettenname) – Nr. 7
%% TODO Anweisung: ein Satz über der ersten Teilaufgabe fehlt in der Bank – Nr. 7
\begin{aufgabe}{Extrempunkte über notwendige und hinreichende Bedingung, Randvergleich auf einem Intervall}
\begin{teile}
\teil $f(x) = x^2 - 6x + 1$. An welcher Stelle liegt der Extrempunkt? x = \leerfeld
\teil $f(x) = x^3 - 12x$. Bestimme die Extremstellen und entscheide mit $f''$, wo ein Maximum liegt. \leerfeld
\end{teile}
\end{aufgabe}

%% TODO Titel: Ich-kann-Satz fehlt in der Bank (Platzhalter: Kettenname) – Nr. 8
%% TODO Anweisung: ein Satz über der ersten Teilaufgabe fehlt in der Bank – Nr. 8
\begin{aufgabe}{Steigung einer Geraden als Quotient aus Höhenunterschied und waagerechtem Unterschied, Steigungsdreieck lesen und zeichnen, Geradengleichung aus Steigung und Punkt (Punkt-Steigungs-Form) – Fehler finden}
\begin{teile}
\teil Für die Gerade durch $A(2 | 3)$ und $B(6 | 11)$ rechnet Ben: \rechnung{m &= 11 - 3 = 8} Finde den Fehler.
\end{teile}
\end{aufgabe}

%% TODO Titel: Ich-kann-Satz fehlt in der Bank (Platzhalter: Kettenname) – Nr. 9
%% TODO Anweisung: ein Satz über der ersten Teilaufgabe fehlt in der Bank – Nr. 9
\begin{aufgabe}{Steigung einer Geraden als Quotient aus Höhenunterschied und waagerechtem Unterschied, Steigungsdreieck lesen und zeichnen, Geradengleichung aus Steigung und Punkt (Punkt-Steigungs-Form) – selbst rechnen}
\begin{teile}
\teil Eine Gerade geht durch $C(1 | 7)$ und $D(5 | -1)$. Berechne ihre Steigung. m = \leerfeld
\end{teile}
\end{aufgabe}
